\section{核心公式与无量纲化}

\subsection{散射电流}

采用时谐约定 $e^{-i\omega t}$，散射电流密度 \cite{Grahn2012}：
\begin{equation}
  \mathbf{J}(\mathbf{r}) = -i\omega\varepsilon_0\bigl(\varepsilon_r(\mathbf{r})-\varepsilon_{r,d}\bigr)\mathbf{E}(\mathbf{r}),
  \label{eq:J}
\end{equation}
其中 $\mathbf{E}$ 为球内总场（孤立球即 Mie 内部场解析解）。
复现中统一去掉公共常数因子 $-i\omega\varepsilon_0$，用
$\widetilde{\mathbf{J}}=(\varepsilon_r-1)\mathbf{E}_{\mathrm{in}}$ 作为积分输入；
物理电流 $J=-i\omega\varepsilon_0\widetilde{\mathbf{J}}$，其绝对尺度在 $C_{\mathrm{sca}}$
归一化到 $\lambda^2/(2\pi)$ 时消去，只有相对空间分布与相位进入多极矩。

\subsection{表2 精确多极矩}

Alaee 表2 的四条公式（讲义 §11，球 Bessel 核，量纲自洽）：

\begin{align}
  \text{ED:}\quad p_\alpha &= -\frac{1}{i\omega}\Bigl\{\int J_\alpha\,j_0(kr)\,d^3r
    + \frac{k^2}{2}\int\bigl[3(\mathbf{r}\cdot\mathbf{J})r_\alpha - r^2J_\alpha\bigr]
    \frac{j_2(kr)}{(kr)^2}\,d^3r\Bigr\}, \label{eq:ed}\\
  \text{MD:}\quad m_\alpha &= \frac32\int(\mathbf{r}\times\mathbf{J})_\alpha\,
    \frac{j_1(kr)}{kr}\,d^3r, \label{eq:md}\\
  \text{EQ:}\quad Q^e_{\alpha\beta} &= -\frac{3}{i\omega}\Bigl\{
    \int\bigl[3(r_\beta J_\alpha + r_\alpha J_\beta) - 2(\mathbf{r}\cdot\mathbf{J})\delta_{\alpha\beta}\bigr]
    \frac{j_1(kr)}{kr}\,d^3r \notag\\
    &\qquad + 2k^2\int\bigl[5r_\alpha r_\beta(\mathbf{r}\cdot\mathbf{J})
    - (r_\alpha J_\beta + r_\beta J_\alpha)r^2 - r^2(\mathbf{r}\cdot\mathbf{J})\delta_{\alpha\beta}\bigr]
    \frac{j_3(kr)}{(kr)^3}\,d^3r\Bigr\}, \label{eq:eq}\\
  \text{MQ:}\quad Q^m_{\alpha\beta} &= 15\int\bigl\{
    r_\alpha(\mathbf{r}\times\mathbf{J})_\beta + r_\beta(\mathbf{r}\times\mathbf{J})_\alpha\bigr\}
    \frac{j_2(kr)}{(kr)^2}\,d^3r. \label{eq:mq}
\end{align}

表1 近似（长波长极限）只是把核常数替换：$1\to j_0$、
$\tfrac13\to\tfrac{j_1}{kr}$、$\tfrac1{15}\to\tfrac{j_2}{(kr)^2}$、
$\tfrac1{105}\to\tfrac{j_3}{(kr)^3}$。因此
\textbf{表1 是表2 的 $kr\to0$ 极限}，两条公式共享同一积分框架，仅核不同。

\subsection{无量纲化（关键设计）}

归一化半径 $u=r/a\in[0,1]$。球内宗量
\begin{equation}
  \rho = k r = (ka)\,u = x\,u,
  \label{eq:rho}
\end{equation}
其中 $k$ 是 \textbf{host 波数}（$k=2\pi/\lambda$），$x=ka=x_{\mathrm{mie}}$。
关键点：\textbf{球 Bessel 辐射核用 host 波数}，内部波数 $k_{\mathrm{in}}=mk$
只通过内部场 $\mathbf{E}_{\mathrm{in}}$ 进入电流空间结构。

位置向量分量写成
\begin{equation}
  \frac{r_x}{a} = u\sin\theta\cos\phi,\quad
  \frac{r_y}{a} = u\sin\theta\sin\phi,\quad
  \frac{r_z}{a} = u\cos\theta.
  \label{eq:rpos}
\end{equation}
\textbf{径向因子 $u$ 不可省略}——这是本轮 blocker 的根因（见第 5 节）。
所有 $r_\alpha$、$\mathbf{r}\cdot\mathbf{J}$、$r^2$ 项都必须含对应的 $u$ 幂次。

\subsection{散射截面（Alaee Eq.1，量纲自洽版）}

从多极矩构造总散射截面 \cite{Alaee2018}，讲义 §12 勘误后的量纲自洽版：
\begin{equation}
  C_{\mathrm{sca}} = \frac{k^4}{6\pi\varepsilon_0^2|E_{\mathrm{inc}}|^2}\Bigl[
    \sum_\alpha\Bigl(|p_\alpha|^2 + \frac{|m_\alpha|^2}{c^2}\Bigr)
    + \frac{1}{120}\sum_{\alpha\beta}\Bigl(
      |kQ^e_{\alpha\beta}|^2 + \frac{|kQ^m_{\alpha\beta}|^2}{c^2}\Bigr) + \cdots
  \Bigr].
  \label{eq:eq1}
\end{equation}
注意：原文印刷 $|m|^2/c$ 是量纲 typo，物理自洽版必须是 $|m|^2/c^2$
与 $|kQ^m|^2/c^2$（讲义 §12 noteworthy）。四极系数 $1/120$ 确认
（$|kQ|^2=k^2|Q|^2$）。

归一化到 $\lambda^2/(2\pi)$（论文 y 轴）。由于 $\lambda^2/(2\pi)=2\pi/k^2$，
从 Eq.~\eqref{eq:eq1} 除以该量后，得到逐多极归一化截面：
\begin{align}
  C_{\mathrm{ED}} &= \frac{x^6}{12\pi^2}\sum_\alpha|p_\alpha|^2, &
  C_{\mathrm{MD}} &= \frac{x^8}{12\pi^2}\sum_\alpha|m_\alpha|^2, \label{eq:c_ed_md}\\
  C_{\mathrm{EQ}} &= \frac{x^8}{1440\pi^2}\sum_{\alpha\beta}|Q^e_{\alpha\beta}|^2, &
  C_{\mathrm{MQ}} &= \frac{x^{10}}{1440\pi^2}\sum_{\alpha\beta}|Q^m_{\alpha\beta}|^2. \label{eq:c_eq_mq}
\end{align}
这些常数的第一性原理推导见第 6.3 节，其中
$1/1440 = 1/(120\times 12)$ 的来源是关键。
